\subsection{PROPERTIES OF THE CARDINALS 0 AND 1}

PROPOSITION 6. Let \((\mathfrak{a}_{\iota})_{\iota \in I}\) be a family of cardinals, and let \(J\) (resp. K) be a subset of \(I\) such that \(\mathfrak{a}_{\iota} = 0\) for all \(\iota \notin J\) (resp. \(\mathfrak{a}_{\iota} = 1\) for all \(\iota \notin K\) ). Then

\[
\sum_ {i \in I} a _ {i} = \sum_ {i \in J} a _ {i} \quad \left(\text { resp. } \underset {i \in I} {\mathbf {P}} a _ {i} = \underset {i \in K} {\mathbf {P}} a _ {i}\right).
\]

The proposition is obvious as regards the sum, for the sum \(S_{I}\) of the family of sets \((a_{i})_{i\in I}\) is equipotent to the union of the sum \(S_{J}\) of the family \((a_{i})_{i\in I}\) and the empty set, and hence equipotent to \(S_{J}\) . The assertion concerning products follows from the fact that the projection \(pr_{K}\) of the product set \(\prod_{i\in I} a_{i}\) onto the partial product \(\prod_{i\in K} a_{i}\) is a bijection (Chapter II, §5, no. 5, Remark 1).

COROLLARY 1. For every cardinal a we have \(a + 0 = a\) . \(1 = a\) .

COROLLARY 2. Let \(\mathfrak{a}\) and \(\mathfrak{b}\) be cardinals and let \(\mathbf{I}\) be a set equipotent to \(\mathfrak{b}\) . For each \(\iota \in \mathbf{I}\) let \(\mathfrak{a}_t = \mathfrak{a}, \mathfrak{c}_t = 1\) . Then

\[
\mathfrak {a}\mathfrak {b} = \sum_ {i \in I} \mathfrak {a} _ {i}, \quad \mathfrak {b} = \sum_ {i \in I} \mathfrak {c} _ {i}.
\]

The second formula is a consequence of the fact that any set is the union of its one-element subsets. The first formula follows from the second by multiplying by \(\mathfrak{a}\) and using Corollary 1.

PROPOSITION 7. Let \((\mathfrak{a}_{i})_{i\in I}\) be a family of cardinals. Then \(\underset{i\in I}{\mathsf{P}}\mathfrak{a}_{i}\neq 0\) if and only if \(\mathfrak{a}_i\neq 0\) for all \(i\in I\) .

This is merely a translation of the condition that a product set should be non-empty (Chapter II, § 5, no. 4, Proposition 5, Corollary 2).

PROPOSITION 8. If \(\mathfrak{a}\) and \(\mathfrak{b}\) are cardinals such that \(\mathfrak{a} + 1 = \mathfrak{b} + 1\) , then \(\mathfrak{a} = \mathfrak{b}\) .

Let \(X = a + 1 = b + 1\) . Then there exist subsets A, B of X with cardinals \(a, b\) , respectively, such that the complements X --- A, X --- B each consist of a single element. Let u, v be these elements. The intersection \(C = A \cap B\) has as complement in X the set \(\{u, v\}\) . If u = v, then A = B = C, so that \(a = b\) . If \(u \neq v\) , then A = C ∪ \(\{v\}\) , B = C ∪ \(\{u\}\) , and therefore \(a = 1 + \text{Card (C)} = b\) .

\pandocbounded{\includegraphics[keepaspectratio]{images/part002_ad0e6982551a0008e23df60b46a8b3d8f866f5ecbe9521df85283d9b8fea8d5c.jpg}}

The reader should beware of assuming that \(a + m = b + m\) implies a = b for all cardinals m (cf.~§6); * we shall see later, however, that this implication is true if m is finite (§5, no. 2, Proposition 3, Corollary 4 and §6, no. 3, Theorem 3, Corollary 4). *

